\chapter{Addition~: premiers pas}\label{ch:g1:addition}

Trois billes rouges et deux billes bleues~: ensemble, combien~?
Mettre ensemble et compter le tout, c'est \emph{ajouter} --- la
premi\`ere op\'eration, avec ses deux signes c\'el\`ebres $+$ et $=$.

\section{Mettre ensemble}

\begin{definition}[Addition]\label{def:g1:addition:def}
\emph{Ajouter}, c'est mettre ensemble et compter le tout. On \'ecrit
\[
3 + 2 = 5
\]
et on lit~: \guillemotleft{} trois plus deux \'egale cinq
\guillemotright{}. Le r\'esultat d'une addition s'appelle la
\emph{somme}\index{somme}.
\end{definition}

\begin{omfigure}
\begin{tikzpicture}[scale=0.6]
  \foreach \i in {0,1,2} \fill[omThm] (\i*1.0,0) circle (0.3);
  \node[font=\large] at (3.0,0) {$+$};
  \foreach \i in {0,1} \fill[omDef] (3.8+\i*1.0,0) circle (0.3);
  \node[font=\large] at (5.8,0) {$=$};
  \foreach \i in {0,1,2} \fill[omThm] (6.6+\i*1.0,0) circle (0.3);
  \foreach \i in {0,1} \fill[omDef] (9.6+\i*1.0,0) circle (0.3);
  \node[below, font=\small] at (5.6,-0.7)
    {$3 + 2 = 5$~: trois et deux font cinq};
\end{tikzpicture}

{\small Une addition en images~: les deux groupes r\'eunis font un
groupe de cinq.}
\end{omfigure}

\begin{example}[L'ordre n'a pas d'importance]\label{ex:g1:addition:order}
$2 + 3$ et $3 + 2$ rassemblent les m\^emes billes~: les deux font
$5$. Pour calculer $2 + 9$, il est plus rapide de partir du plus grand
nombre~: dire \guillemotleft{} neuf \guillemotright{}, puis compter
deux de plus --- \guillemotleft{} dix, onze \guillemotright{}. Donc
$2 + 9 = 11$.
\end{example}

\section{Les amis de 10}

\begin{example}[Compl\'ements \`a 10]\label{ex:g1:addition:friends}
Deux nombres qui font $10$ ensemble sont des \guillemotleft{} amis de
$10$ \guillemotright{}~:
\[
1 + 9, \quad 2 + 8, \quad 3 + 7, \quad 4 + 6, \quad 5 + 5 .
\]
Les dix doigts les montrent tous~: lever $3$ doigts, et les $7$
pli\'es sont l'ami. Conna\^itre les amis de $10$ par c{\oe}ur rend
beaucoup d'additions rapides.
\end{example}

\begin{omfigure}
\begin{tikzpicture}[scale=0.72]
  % ten-frame: 2 rows of 5 cells
  \draw[thick] (0,0) grid[step=1] (5,2);
  \draw[very thick] (0,0) rectangle (5,2);
  \foreach \i in {0,1,2} \fill[omDef] (\i+0.5,1.5) circle (0.3);
  \foreach \i in {3,4} \draw[omThm, thick] (\i+0.5,1.5) circle (0.3);
  \foreach \i in {0,...,4} \draw[omThm, thick] (\i+0.5,0.5) circle (0.3);
  \node[right, font=\small, align=left] at (5.6,1.0)
    {un \emph{cadre de 10}~: $10$ cases.\\
     $3$ remplies, $7$ vides~: $3 + 7 = 10$};
\end{tikzpicture}

{\small Le cadre de $10$ montre d'un coup d'{\oe}il chaque paire
d'amis de $10$~: ce qui est rempli et ce qui manque font toujours
$10$ ensemble.}
\end{omfigure}

\begin{method}[Passer par 10]\label{met:g1:addition:crossten}
Pour calculer $8 + 5$~:
\begin{enumerate}
  \item prendre dans le $5$ ce qu'il manque au $8$ pour atteindre
        $10$~: c'est $2$ (son ami de $10$)~;
  \item $8 + 2 = 10$, et il reste $3$ du $5$~;
  \item $10 + 3 = 13$. Donc $8 + 5 = 13$.
\end{enumerate}
Passer par $10$ transforme une addition difficile en deux faciles.
\end{method}

\begin{omfigure}
\begin{tikzpicture}[scale=0.62]
  % ten-frame with 8 blue, 2 red (borrowed from the 5)
  \draw[thick] (0,0) grid[step=1] (5,2);
  \draw[very thick] (0,0) rectangle (5,2);
  \foreach \i in {0,...,4} \fill[omDef] (\i+0.5,1.5) circle (0.3);
  \foreach \i in {0,1,2} \fill[omDef] (\i+0.5,0.5) circle (0.3);
  \foreach \i in {3,4} \fill[omThm] (\i+0.5,0.5) circle (0.3);
  % 3 red left outside the frame
  \foreach \i in {0,1,2} \fill[omThm] (6.2+\i,1.0) circle (0.3);
  \node[below, font=\small] at (2.5,-0.25) {le cadre est plein~: $10$};
  \node[below, font=\small] at (7.2,0.55) {$3$ d\'ebordent};
\end{tikzpicture}

{\small $8 + 5$ sur un cadre de $10$~: les $8$ jetons bleus prennent
$2$ des $5$ rouges pour remplir le cadre --- une dizaine pleine et
$3$ de plus, donc $13$.}
\end{omfigure}

\begin{omfigure}
\begin{tikzpicture}[scale=0.6]
  \draw[-Stealth] (-0.4,0) -- (15.4,0);
  \foreach \i in {0,...,15} \draw (\i,-0.1) -- (\i,0.1);
  \foreach \i in {0,5,8,10,13,15}
    \node[below, font=\footnotesize] at (\i,-0.12) {\i};
  \draw[-Stealth, omDef, thick] (8,0.3) .. controls (9,0.95) ..
    (10,0.3);
  \node[omDef, font=\small] at (9,1.05) {$+2$};
  \draw[-Stealth, omThm, thick] (10,0.3) .. controls (11.5,1.15) ..
    (13,0.3);
  \node[omThm, font=\small] at (11.5,1.2) {$+3$};
\end{tikzpicture}

{\small $8 + 5$ sur la droite num\'erique~: d'abord un saut jusqu'\`a
$10$, puis le reste.}
\end{omfigure}

\begin{example}[Les doubles]\label{ex:g1:addition:doubles}
Les doubles m\'eritent d'\^etre appris par c{\oe}ur~:
\[
1+1=2, \quad 2+2=4, \quad 3+3=6, \quad 4+4=8, \quad 5+5=10,
\]
\[
6+6=12, \quad 7+7=14, \quad 8+8=16, \quad 9+9=18, \quad 10+10=20 .
\]
Les presque-doubles viennent gratuitement~: $6 + 7$ c'est un de plus
que $6 + 6$, donc $13$.
\end{example}

\section{Petits probl\`emes}

\begin{example}\label{ex:g1:addition:problem}
\guillemotleft{} L\'ea a $7$ autocollants. Elle en re\c{c}oit $6$ de
plus. Combien en a-t-elle maintenant~? \guillemotright{} Recevoir
plus, c'est ajouter~: $7 + 6 = 13$ (en passant par $10$~:
$7 + 3 = 10$, puis $+3$). L\'ea a $13$ autocollants. Toujours finir
par une petite phrase~!
\end{example}

\section{Exercices}

\begin{exercise}[$\star$]\label{exo:g1:addition:1}
Dessine les billes et calcule~: $4 + 3$~; \quad $5 + 2$~; \quad
$6 + 4$.
\end{exercise}

\begin{exercise}[$\star$]\label{exo:g1:addition:2}
Calcule en partant du plus grand nombre~: $3 + 8$~; \quad
$2 + 12$~; \quad $4 + 9$.
\end{exercise}

\begin{exercise}[$\star$]\label{exo:g1:addition:3}
Donne l'ami de $10$ de~: $7$~; \ $4$~; \ $9$~; \ $5$~; \ $1$.
\end{exercise}

\begin{exercise}[$\star$]\label{exo:g1:addition:4}
Calcule en passant par $10$ (\cref{met:g1:addition:crossten})~:
$9 + 4$~; \quad $8 + 7$~; \quad $6 + 5$.
\end{exercise}

\begin{exercise}[$\star$]\label{exo:g1:addition:5}
R\'ecite les doubles, puis calcule en s'en servant~: $7 + 8$~; \quad
$5 + 6$~; \quad $9 + 8$.
\end{exercise}

\begin{exercise}[$\star$]\label{exo:g1:addition:6}
Compl\`ete les trous~:
\[
6 + \;?\; = 10, \qquad
\;?\; + 5 = 12, \qquad
10 + \;?\; = 17 .
\]
\end{exercise}

\begin{exercise}[$\star$]\label{exo:g1:addition:7}
\guillemotleft{} Dans le bus, il y a $8$ enfants~; $5$ adultes
montent. Combien de personnes y a-t-il dans le bus~?
\guillemotright{} \'Ecris l'addition et la phrase r\'eponse.
\end{exercise}

\begin{exercise}[$\star$]\label{exo:g1:addition:8}
Tom a $6$ voitures rouges et $7$ voitures bleues. Mia a $10$
voitures. Qui a le plus de voitures~?
\end{exercise}

\begin{exercise}[$\star\star$]\label{exo:g1:addition:9}
Trois d\'es montrent $4$, $6$ et $5$. Quel est le total~? (Ajoute
d'abord deux d\'es --- choisis les deux qui forment un ami de $10$~!)
\end{exercise}
